\customchapter{RESUMEN}

Los modelos de aprendizaje automático que se usan en producción pierden calidad cuando los
datos que reciben dejan de parecerse a los datos con los que fueron entrenados. En un
\textit{pipeline} de dos etapas, formado por una etapa de preparación de datos y un modelo
predictivo, este problema puede originarse en cualquiera de las dos etapas. Sin embargo, la
respuesta habitual es volver a entrenar todo el \textit{pipeline}, lo que resulta costoso y
no siempre es necesario. Esta tesis propone DARL (\textit{Drift-Aware Reinforcement
Learner}), un enfoque que trata el mantenimiento del \textit{pipeline} como una secuencia de
decisiones. Cada día, un agente de aprendizaje por refuerzo revisa señales de monitoreo y
elige entre no hacer nada, corregir solo la preparación de los datos, volver a entrenar solo
el modelo o volver a entrenar todo el \textit{pipeline}. Como el agente no observa
directamente el tipo de cambio que ocurre, el problema se formula como un Proceso de
Decisión de Markov Parcialmente Observable y se resuelve con una red Q profunda (DQN). La
recompensa combina la mejora en el área bajo la curva de precisión y exhaustividad (AUPRC)
con el costo de cada acción, y se calcula con un día de retraso, cuando se conocen las
etiquetas reales. El caso de estudio es la predicción temprana de sepsis con los datos del
PhysioNet/Computing in Cardiology Challenge 2019: un \textit{pipeline} entrenado con
pacientes de un hospital se usa en una unidad de cuidados intensivos simulada de 200 camas
de otro hospital, donde se inyectan de forma controlada cambios en la medición de los
signos vitales, cambios en el etiquetado de la sepsis y su combinación. La política
aprendida se contrasta con reglas de referencia, como no actuar nunca, volver a entrenar
siempre, actuar al azar o seguir un umbral fijo.

\noindent \textbf{Palabras clave:}\\
\noindent Cambio de distribución; Aprendizaje por refuerzo; Mantenimiento de modelos; Predicción de sepsis
